% R597: whole EGA I Chapter0 3.5.2–3.
% OpenAI Codex, GPT-6.1 Sol, Ultra; no human review claimed.
\begin{env}[3.5.2]
\label{0.3.5.2-fr}
Bajo las hipótesis de (3.4.3), sean $\mathcal{F}$, $\mathcal{G}$,
$\mathcal{H}$ prehaces con valores en $\boldsymbol{K}$ sobre $X$, $Y$, $Z$
respectivamente, y sean
$u:\mathcal{G}\to\psi_*(\mathcal{F})$,
$v:\mathcal{H}\to\psi'_*(\mathcal{G})$ un $\psi$-morfismo y un
$\psi'$-morfismo respectivamente. Se deduce de ellos un $\psi''$-morfismo
\[
  w:\mathcal{H}\xrightarrow{v}\psi'_*(\mathcal{G})
  \xrightarrow{\psi'_*(u)}
  \psi'_*(\psi_*(\mathcal{F}))=\psi''_*(\mathcal{F}),
\]
que se llama, por definición, el \emph{compuesto} de $u$ y de $v$. Se puede,
pues, considerar los pares $(X,\mathcal{F})$ formados por un espacio topológico
$X$ y un prehaz $\mathcal{F}$ sobre $X$ (con valores en $\boldsymbol{K}$)
como una \emph{categoría}, cuyos morfismos son los pares
$(\psi,\theta):(X,\mathcal{F})\to(Y,\mathcal{G})$ formados por una aplicación
continua $\psi:X\to Y$ y un $\psi$-morfismo
$\theta:\mathcal{G}\to\mathcal{F}$.
\end{env}

\begin{env}[3.5.3]
\label{0.3.5.3-fr}
Sean $\psi:X\to Y$ una aplicación continua y $\mathcal{G}$ un
\emph{prehaz} sobre $Y$ con valores en $\boldsymbol{K}$. Llamaremos
\emph{imagen inversa de $\mathcal{G}$ por $\psi$} a un par
$(\mathcal{G}',\rho)$, donde $\mathcal{G}'$ es un \emph{haz} sobre $X$ con
valores en $\boldsymbol{K}$, y
$\rho:\mathcal{G}\to\mathcal{G}'$ es un $\psi$-morfismo (es decir, un
homomorfismo $\mathcal{G}\to\psi_*(\mathcal{G}'))$ tales que, para todo
\emph{haz} $\mathcal{F}$ sobre $X$ con valores en $\boldsymbol{K}$,
la aplicación
\begin{equation*}
  \operatorname{Hom}_X(\mathcal{G}',\mathcal{F})
  \longrightarrow \operatorname{Hom}_\psi(\mathcal{G},\mathcal{F})
  =\operatorname{Hom}_Y(\mathcal{G},\psi_*(\mathcal{F}))
\tag{3.5.3.1}\label{0.3.5.3.1-fr}
\end{equation*}
que transforma $v$ en $\psi_*(v)\circ\rho$ sea una \emph{biyección}~; esta
aplicación, al ser funtorial en $\mathcal{F}$, define entonces un
isomorfismo de funtores en $\mathcal{F}$. Al ser el par
$(\mathcal{G}',\rho)$ solución de un problema universal, se sabe que está
\emph{determinado salvo un isomorfismo único} cuando existe. Se escribirá
entonces $\mathcal{G}'=\psi^*(\mathcal{G})$, $\rho=\rho_{\mathcal{G}}$, y por
abuso de lenguaje se dirá que $\psi^*(\mathcal{G})$ es \emph{el haz imagen
inversa} de $\mathcal{G}$ por $\psi$, entendiéndose que
$\psi^*(\mathcal{G})$ se considera provisto del $\psi$-\emph{morfismo
canónico}
$\rho_{\mathcal{G}}:\mathcal{G}\to\psi^*(\mathcal{G})$, es decir, del
\emph{homomorfismo canónico} de prehaces sobre $Y$~:
\begin{equation*}
  \rho_{\mathcal{G}}:\mathcal{G}\to\psi_*(\psi^*(\mathcal{G})).
\tag{3.5.3.2}\label{0.3.5.3.2-fr}
\end{equation*}
Para todo homomorfismo
$v:\psi^*(\mathcal{G})\to\mathcal{F}$ (donde $\mathcal{F}$ es un haz sobre
$X$ con valores en $\boldsymbol{K}$), se pondrá
$v^\flat=\psi_*(v)\circ\rho_{\mathcal{G}}:
\mathcal{G}\to\psi_*(\mathcal{F})$. Por definición, \emph{todo} morfismo de
prehaces $u:\mathcal{G}\to\psi_*(\mathcal{F})$ es de la forma $v^\flat$ para
un único $v$, que se denotará por $u^\sharp$. En otras palabras, todo
morfismo $u:\mathcal{G}\to\psi_*(\mathcal{F})$ de prehaces se factoriza
de manera única como
\begin{equation*}
  u:\mathcal{G}\xrightarrow{\rho_{\mathcal{G}}}
  \psi_*(\psi^*(\mathcal{G}))
  \xrightarrow{\psi_*(u^\sharp)}\psi_*(\mathcal{F}).
\tag{3.5.3.3}\label{0.3.5.3.3-fr}
\end{equation*}
\end{env}
